\section{Finite element formulation}\label{sec:fe}

This section states the continuum equations and the discrete operators that \fea{} assembles. The formulation is the standard displacement-based
one \cite{zienkiewicz2013fem,bathe2014fep}; we state it to fix notation and to make clear which parts of the code correspond to which equation.

\subsection{Weak form and discretisation}
For a body $\Omega$ with displacement field $\bm u(\bm x,t)$, density $\rho$, Cauchy stress $\bm\sigma$, body force $\bm b$ and traction $\bm t$ on $\Gamma_t$,
the principle of virtual work reads
\begin{equation}
  \int_\Omega \rho\,\delta\bm u\cdot\ddot{\bm u}\,\mathrm d\Omega
  +\int_\Omega \delta\bm\varepsilon:\bm\sigma\,\mathrm d\Omega
  =\int_\Omega\delta\bm u\cdot\bm b\,\mathrm d\Omega+\int_{\Gamma_t}\delta\bm u\cdot\bm t\,\mathrm d\Gamma
  \quad\forall\,\delta\bm u .
  \label{eq:weak}
\end{equation}
Approximating $\bm u\approx\bm N\vu$ with shape functions $\bm N$ and $\bm\varepsilon=\bm B\vu$, and taking linear elasticity $\bm\sigma=\bm D\bm\varepsilon$, gives the
semi-discrete equations
\begin{equation}
  \M\ddot{\vu}+\C\dot{\vu}+\K\vu=\F,
  \label{eq:sd}
\end{equation}
with $\K=\sum_e\K_e$ and $\M=\sum_e\M_e$ assembled from the element matrices \eqref{eq:elem}.
The damping matrix $\C$ is either Rayleigh damping $\C=\alpha\M+\beta\K$, with $(\alpha,\beta)$ calibrated to a target ratio $\zeta$ at two angular frequencies
$\omega_i,\omega_j$ by solving $\tfrac12\bigl(\alpha/\omega_k+\beta\omega_k\bigr)=\zeta$, $k\in\{i,j\}$ (\texttt{RayleighDamping.calibrate}),
or a modal damping ratio applied directly to decoupled modal equations (\texttt{ModalDamping}).
A lumped (diagonal) mass matrix is also available for explicit integration.

\subsection{Isoparametric elements and quadrature}
Two- and three-dimensional elements use isoparametric mapping: the geometry is interpolated with the same shape functions as the displacement,
$\bm x=\sum_a N_a(\bm\xi)\bm x_a$. With Jacobian $\bm J=\partial\bm x/\partial\bm\xi$, Cartesian derivatives are
$\partial N_a/\partial\bm x=\bm J^{-\top}\partial N_a/\partial\bm\xi$, and integrals over the element are evaluated by Gauss quadrature,
\begin{equation}
  \K_e\approx\sum_{p}w_p\,\bm B^{\!\top}(\bm\xi_p)\,\bm D\,\bm B(\bm\xi_p)\,\det\bm J(\bm\xi_p)\,t,
  \label{eq:gauss}
\end{equation}
where $t$ is the thickness for plane problems. Strains use engineering shear and the Voigt order $[xx,yy,xy]$ in 2-D and $[xx,yy,zz,xy,yz,xz]$ in 3-D.

\subsection{Element families}
\Cref{tab:elements} lists the element families. The registry holds \nElements{} entries; variants (for example selective-reduced and full integration of the same element) are
counted separately.

\begin{table}[t]
\centering\footnotesize
\caption{Element families in \fea{}. Class names are those in \texttt{fea\_engine.elements}.}
\label{tab:elements}
\begin{tabular}{>{\raggedright\arraybackslash}p{3.0cm}>{\raggedright\arraybackslash}p{5.2cm}>{\raggedright\arraybackslash}p{5.2cm}}
\toprule
Family & Classes & Formulation notes \\
\midrule
Trusses & \texttt{TrussTL2D}, \texttt{TrussPlastic2D} & total-Lagrangian, Green--Lagrange strain; 1-D elastoplastic variant \\
Beams, 2-D & \texttt{Beam2DEulerBernoulli}, \texttt{Beam2DCorotational}, \texttt{Beam2DReissner} & Hermite beam; corotational large rotation; shear-deformable planar beam \\
Beams, 3-D & \texttt{Beam3DEulerBernoulli}, \texttt{Beam3DCorotational} & per-element local axes; corotational variant \\
Plane stress/strain & \texttt{Tri3}, \texttt{Tri6}, \texttt{Quad4}, \texttt{Quad8}\texttt{PlaneStress} & isoparametric, Gauss quadrature \\
Solids & \texttt{Tet4}, \texttt{Tet10}, \texttt{Hex8}, \texttt{Hex20}\texttt{Solid3D}, \texttt{Hex8SolidBbar} & B-bar variant against volumetric locking \\
Plate & \texttt{Quad4MindlinPlate} & Mindlin--Reissner, selective reduced integration \\
Shells & \texttt{Shell4MITC}, \texttt{Shell4MITCCorotational}, \texttt{Shell4Director} & MITC4 assumed shear strain; corotational; degenerated-shell director kinematics (linear stiffness only, checked against MITC4) \\
Inelastic/hyperelastic & \texttt{Hex8PlasticJ2}, \texttt{Hex8PlasticJ2Kinematic}, \texttt{Quad4PlasticJ2PlaneStress}, \texttt{Tet4NeoHookean}, \texttt{Tet10SolidTL} & J2 plasticity; compressible Neo-Hookean; St~Venant--Kirchhoff-type total-Lagrangian 10-node tetrahedron \\
Contact & \texttt{GapContactPenalty}, \texttt{GapContactCurvedFriction}, \texttt{NodeToSegment\allowbreak Contact2D} (and friction variant) & penalty gap; node-to-segment with Coulomb friction \\
\bottomrule
\end{tabular}
\end{table}

\paragraph{Locking remedies.}
Low-order elements lock in the near-incompressible and thin limits. For the 8-node hexahedron, \texttt{Hex8SolidBbar} replaces the volumetric part of $\bm B$ by its element average,
\begin{equation}
  \bar{\bm B}_p=\bm B_p-\tfrac13\bm m\bm m^{\!\top}\bm B_p+\bar{\bm B}_{\mathrm{vol}},\qquad
  \bar{\bm B}_{\mathrm{vol}}=\frac{1}{V_e}\sum_q w_q\det\bm J_q\;\tfrac13\bm m\bm m^{\!\top}\bm B_q,
  \label{eq:bbar}
\end{equation}
with $\bm m=[1,1,1,0,0,0]^{\!\top}$ \cite{hughes1980bbar}. The 4-node Mindlin plate uses selective reduced integration of the shear term.
\texttt{Shell4MITC} follows the MITC4 element \cite{dvorkin1984}: the covariant transverse shear strains are sampled at the four edge mid-points and interpolated,
\begin{equation}
  \tilde e_{rt}=\tfrac12(1+s)\,e_{rt}^{(A)}+\tfrac12(1-s)\,e_{rt}^{(B)},\qquad
  \tilde e_{st}=\tfrac12(1+r)\,e_{st}^{(C)}+\tfrac12(1-r)\,e_{st}^{(D)},
  \label{eq:mitc}
\end{equation}
which removes shear locking in thin-shell bending. Membrane parts use the plane-stress $\bm B$ and bending parts the Mindlin curvature matrix. The element has a small penalty stiffness on the drilling rotation to
avoid singular matrices on flat patches. As in classic MITC4, membrane locking in coarse curved meshes is not addressed, and each element is a flat facet built from its own corner nodes.

\subsection{Assembly}
\Cref{alg:assembly} describes assembly. For a vectorised element the loop over $e$ is replaced by a batched kernel; otherwise the loop calls the element methods. Coordinates and
element constants may vary in space through coefficient fields (\cref{sec:recovery}).

\begin{algorithm}[t]
\caption{Global assembly of $\K$ (and, identically, $\M$)}\label{alg:assembly}
\begin{algorithmic}[1]
\Require mesh with connectivity $\{\mathcal I_e\}$, element formulation, constitutive matrix $\bm D$
\State $\K\gets\bm 0$ \Comment{sparse or dense, set at construction}
\For{each element $e$ of each block}
  \State $\bm x_e\gets$ nodal coordinates of $e$
  \State $\K_e\gets\sum_p w_p\bm B_p^{\!\top}\bm D\bm B_p\det\bm J_p\,t$ \Comment{element \texttt{stiffness}}
  \State $\mathcal G_e\gets$ global DOF indices from $\mathcal I_e$ and $n_{\mathrm{dof/node}}$
  \State $\K[\mathcal G_e,\mathcal G_e]\mathrel{+}=\K_e$
\EndFor
\State add spring and foundation terms (\cref{sec:constraints})
\end{algorithmic}
\end{algorithm}
